% Compile this UTF-8 document with XeLaTeX or LuaLaTeX.
\documentclass[11pt]{article}
\usepackage{fontspec}
\usepackage{amsmath,amssymb,mathtools,bm}
\usepackage[margin=1in]{geometry}
\setlength{\parindent}{0pt}
\setlength{\parskip}{0.5em}
\begin{document}
\section*{1.1 What Is Regression Analysis?}
{\small Source: Regression Modeling with Actuarial and Financial Applications (Frees)/Regression Modeling with Actuarial and Financial Applications.pdf\par}
\medskip
1
 Regression and the Normal Distribution
Chapter Preview. Regression analysis is a statistical method that is widely used in many
fields of study, with actuarial science being no exception. This chapter provides an intro-
 duction to the role of the normal distribution in regression, the use of logarithmic trans-
formations in specifying regression relationships, and the sampling basis that is critical
for inferring regression results to broad populations of interest.
1.1 What Is Regression Analysis?
Statistics is about data. As a discipline, it is about the collection, summarization,
and analysis of data to make statements about the real world. When analysts
 collect data, they are really collecting information that is quantified, that is, trans-
 formed to a numerical scale. There are easy, well-understood rules for reducing
 the data, through either numerical or graphical summary measures. These sum-
mary measures can then be linked to a theoretical representation, or model, of
the data. With a model that is calibrated by data, statements about the world can
 be made. Statistics is about
 Statistical methods have had a major impact on several fields of study: the collection,
 summarization, and
: In the area of data collection, the careful design of sample surveys is crucial analysis of data to
to market research groups and to the auditing procedures of accounting  make staremens
 about the real world.
firms.
: Experimental design is a subdiscipline devoted to data collection. The focus
of experimental design is on constructing methods of data collection that
will extract information in the most efficient way possible. This is espe-
cially important in fields such as agriculture and engineering where each
observation is expensive, possibly costing millions of dollars.
- Other applied statistical methods focus on managing and predicting data.
Process control deals with monitoring a process over time and deciding
 when intervention is most fruitful. Process control helps manage the quality
 of goods produced by manufacturers.
- Forecasting is about extrapolating a process into the future, whether it be
 sales of a product or movements of an interest rate.
Regression analysis is a statistical method used to analyze data. As we will see,
 the distinguishing feature of this method is the ability to make statements about
1
\clearpage
2 Regression and the Normal Distribution
Table 1.1 Galton's
1885 Regression Data Height of Parents' Height
Adult Child
in Inches \(< 6 4 . 0\) 64.5 65.5 66.5 67.5 68.5 69.5 70.5 71.5 72.5 >73.0 Total
>73.7 5324 14
73.2 3 \(4\) 322 3 17
72.2 1 4411497 1 41
71.2 2 111820742 64
70.2 541921251410 1 99
69.2 1 271338 4833 18 5 2 167
68.2 1 71428 3420123 1 120
67.2 2 51117 3831273 4 138
66.2 2 5111736251713 117
65.2 1 17 2151641 1 48
64.2 4 4 55141116 59
63.2 2 493 5711 32
62.2 1 3 3 7
<61.2 1 1 1 1 1 5
Total 14 23 66 78 211 219 183 68 43 19 4 928
Source: Stigler (1986).
variables after having controlled for values of known explanatory variables. As
important as other methods are, it is regression analysis that has been the most
 influential one. To illustrate, an index of business journals, ABI/INFORM, lists
 more than 24,000 articles using regression techniques over the thirty-year period
1978-- 2007.And these are only the applications that were considered innovative
enough to be published in scholarly reviews!
 Regression analysis of data is so pervasive in modern business that it is easy
to overlook the fact that the methodology is barely more than 120 years old.
Scholars attribute the birth of regression to the 1885 presidential address of Sir
Francis Galton to the anthropological section of the British Association of the
Advancement of Sciences. In that address, described in Stigler (1986), Galton
provided a description of regression and linked it to normal curve theory. His
@ EMPIRICAL discovery arose from his studies of properties of natural selection and inheritance.
 Filename is To illustrate a dataset that can be analyzed using regression methods, Table 1.1
"Galton" displays some data included in Galton's 1885 paper. The table displays the heights
of 928 adult children, classified by an index of their parents' height. Here, all
female heights were multiplied by \(1 . 0 8\) , and the index was created by taking the
 average of the father's height and rescaled mother's height. Galton was aware that
the parents' and the adult child's height could each be adequately approximated
Regression analysis is by a normal curve. In developing regression analysis, he provided a single model
 a method to quantify for the joint distribution of heights.
 the relationship Table 1.1 shows that much of the information concerning the height of an adult
 between a variable of
interest and child can be attributed to, or " plained," in terms of the parents' height. Thus,
 explanatory variables. we use the term explanatory variable for measurements that provide information
\clearpage

\end{document}
% Compile this UTF-8 document with XeLaTeX or LuaLaTeX.
\documentclass[11pt]{article}
\usepackage{fontspec}
\usepackage{amsmath,amssymb,mathtools,bm}
\usepackage[margin=1in]{geometry}
\setlength{\parindent}{0pt}
\setlength{\parskip}{0.5em}
\begin{document}
\section*{1.2 Fitting Data to a Normal Distribution}
{\small Source: Regression Modeling with Actuarial and Financial Applications (Frees)/Regression Modeling with Actuarial and Financial Applications.pdf\par}
\medskip
1.2 Fitting Data to a Normal Distribution 3
Figure 1.1 Ten
 deutsche mark --
AY2252 featuring the scientist
German currency
Gauss and the normal
curve.
AY5229324U0
 about a variable of interest. Regression analysis is a method to quantify the rela-
 tionship between a variable of interest and explanatory variables. The methodol-
 ogy used to study the data in Table 1.1 can also be used to study actuarial and
other risk management problems, the thesis of this book.
1.2 Fitting Data to a Normal Distribution
Historically, the normal distribution had a pivotal role in the development of
regression analysis. It continues to play an important role, although we will be
 interested in extending regression ideas to highly "- nonnormal"data.
Formally, the normal curve is defined by the function
\[
\mathrm{f} ( y )=\frac{1} {\sigma\, \sqrt{2 \pi}} \operatorname{e x p} \left(-\frac{1} {2 \sigma^{2}} \left( y-\mu\right)^{2} \right) . \tag{1.1}
\]
 This curve is a probability density function with the whole real line as its domain. Appendix A3.1
 parameters, \(\mu{}\) and \(\sigma^{2}\) (the mean and provides additional
From equation (1.1), we see that the curve is symmetric about \(\mu{}\) \(\sigma^{2}\) . These two  normal curve.
 median). The degree of peakedness is controlled by the parameter details about the
 are known as the location and scale parameters, respec- including a graph and
tively. Appendix A3.1 provides additional details about this curve, including a distribution table.
 graph and tables of its cumulative distribution that we will use throughout the
text.
The normal curve is also depicted in Figure 1.1, a display of an out-of-date
German currency note, the ten Deutsche Mark. This note contains the image
of the German Carl Gauss, an eminent mathematician whose name is often
linked with the normal curve (it is sometimes referred to as the Gaussian curve).
Gauss developed the normal curve in connection with the theory of least squares
for fitting curves to data in 1809, about the same time as related work by the
French scientist Pierre LaPlace. According to Stigler (1986), there was quite a
bit of acrimony between these two scientists about the priority of discovery! The
normal curve was first used as an approximation to histograms of data around
1835 by Adolph Quetelet, a Belgian mathematician and social scientist. As with
 many good things, the normal curve had been around for some time, since about
1720, when Abraham de Moivre derived it for his work on modeling games of
\clearpage
4 Regression and the Normal Distribution
Table 1.2 Summary
Statistics of Standard 25th 75th
Massachusetts Variable Number Mean Median Deviation Minimum Maximum Percentile Percentile
Automobile Bodily
 Injury Claims Claims2720.481 0.793 1.101 -3.101 3.912 -0.114 1.168
Note: Data are in logs of thousands of dollars.
chance. The normal curve is popular because it is easy to use and has proved
 successful in many applications.
@ EMPIRICAL
Filename is Example: Massachusetts Bodily Injury Claims. For our first look at fitting
"MassBodilyInjury" the normal curve to a set of data, we consider data from Rempala and Derrig
(2005). They considered claims arising from automobile bodily injury insurance
coverages. These are amounts incurred for outpatient medical treatments that
arise from automobile accidents, typically sprains, broken collarbones, and the
like. The data consist of a sample of 272 claims from Massachusetts that were
closed in 2001 (by " closed, we mean that the claim is settled and no additional
liabilities can arise from the same accident). Rempala and Derrig were interested
in developing procedures for handling mixtures of " typical"claims and others
from providers who reported claims fraudulently. For this sample, we consider
 only those typical claims, ignoring the potentially fraudulent ones.
Table 1.2 provides several statistics that summarize different aspects of the
distribution. Claim amounts are in units of logarithms of thousands of dollars.
The average logarithmic claim is 0.481, corresponding to \$1,617.77 (=1000
exp(0.481)). The smallest and largest claims are -3.101 (\$45) and 3.912
(\$50,000), respectively.
For completeness, here are a few definitions. The sample is the set of data
 available for analysis, denoted by \(y_{1} , \dots, \, y_{n}\) . Here, \(n\) is the number of observa-
tions, \(y_{1}\) represents the first observation, \(y_{2}\) the second, and so on up to \(y_{n}\) for the
\(n\) th observation. Here are a few important summary statistics.
Basic Summary Statistics
(i) The mean is the average of observations, that is, the sum of the observations
divided by the number of units. Using algebraic notation, the mean is
\[
\begin{aligned} {\overline{{y}}} & {{}=\frac{1} {n} \left( y_{1}+\cdots+y_{n} \right)=\frac{1} {n} \sum_{i=1}^{n} y_{i} .} \\ \end{aligned}
\]
(ii) The median is the middle observation when the observations are ordered
by size. That is, it is the observation at which 50\% are below it (and 50\%
are above it).
\clearpage
1.2 Fitting Data to a Normal Distribution 5
Density outlier Figure 1.2 Bodily
0.3 injury relative
frequency with
0.2 normal curve
0.1 median 75th percentlesuperimposed.
25th percentilel
0.0 2 0 2 \(\overline{{8}}\) Toutliers Eh
of bodily injury
LOGCLAIMS
l(ii) The standard deviation is a measure of the spread, or scale, of the
distribution. It is computed as
\[
\begin{aligned} {s_{y}} & {{}=\sqrt{\frac{1} {n-1} \sum_{i=1}^{n} \left( y_{i} \,-\, \overline{{y}} \right)^{2}} .} \\ \end{aligned}
\]
(iv) A percentile is a number at which a specified fraction of the observations
is below it, when the observations are ordered by size. For example, the
 25th percentile is the number so that 25\% of observations are below it.
To help visualize the distribution, Figure 1.2 displays a histogram of the
data. Here, the height of the each rectangle shows the relative frequency of
 observations that fall within the range given by its base. The histogram provides
 a quick visual impression of the distribution; it shows that the range of the data is
approximately \((-4 , 4 )\) that the central tendency is slightly greater than zero, and
 that the distribution is roughly symmetric.
 Normal Curve Approximation
Figure 1.2 also shows a normal curve superimposed, using \(\overline{{y}}\) for \(\mu{}\) and \(s_{y}^{2}\) for \(\sigma^{2}\) .
With the normal curve, only two quantities ( \([ \mu\mathrm{~ a n d ~} \sigma^{2} ]\) are required to summarize
the entire distribution. For example, Table 1.2 shows that 1.168 is the 75th
percentile, which is approximately the \(2 0 4 \mathrm{t h} \ (=. 7 5 \times2 7 2 )\) largest observation
from the entire sample. From the equation (1.1) normal distribution, we see that
\(z=( y-\mu) / \sigma\) is a standard normal, of which 0.675 is the 75th percentile. Thus,
\(\overline{y} \,+\, 0 . 6 7 5 s_{y} \,=\, 0 . 4 8 1 \,+\, 0 . 6 7 5 \, \times\, 1 . 1 0 1 \,=\, 1 . 2 2 4 ~ \mathrm{i s}\) the 75th percentile using the
 normal curve approximation.
Box Plot
 A quick visual inspection of a variable's distribution can reveal some surprising
 features that are hidden by statistics: numerical summary measures. The box plot,
also known as a box-and-whiskers plot, is one such graphical device. Figure 1.3
illustrates a box plot for the bodily injury claims. Here, the box captures the
\clearpage
6 Regression and the Normal Distribution
Figure 1.4 Redraw- Density Sample Quantiles
ing of Figure 1.2 with 0.7
 an increased number 0.6 3
of rectangles. 0.5 2
Figure 1.5 A \(q q\) plot 0.4 1
of bodily injury 0.3 0
claims, using a normal 0.2 1
reference distribution.
0.1 2
0.0 3
3 2 1 0 1 2 3 3 2 1 0 1 2 3
LOGCLAIMS Theoretical Quantiles
middle 50\% of the data, with the three horizontal lines corresponding to the
75th, 50th, and 25th percentiles, reading from top to bottom. The horizontal lines
above and below the box are the " whislers." The upper whisker is \(1 . 5\) times
the interquartile range (the difference between the 75th and 25th percentiles)
 above the 75th percentile. Similarly, the lower whisker is 1.5 times the interquar-
 tile range below the 25th percentile. Individual observations outside the whiskers
are denoted by small circular plotting symbols and are referred to as " outliers?
Graphs are powerful tools; they allow analysts to readily visualize nonlin-
ear relationships that are hard to comprehend when expressed verbally or by
mathematical formula. However, by their very flexibility, graphs can also readily
 deceive the analyst. Chapter 21 will underscore this point. For example, Fig-
ure \(1 . 4\) is a redrawing of Figure 1.2; the difference is that Figure 1.4 uses more,
and finer, rectangles. This finer analysis reveals the asymmetric nature of the
 sample distribution that was not evident in Figure 1.2.
 Quantile-Quantile Plots
Increasing the number of rectangles can unmask features that were not previ-
 ously apparent; however, there are, in general, fewer observations per rectangle,
meaning that the uncertainty of the relative frequency estimate increases. This
 represents a trade-off. Instead of forcing the analyst to make an arbitrary decision
about the number of rectangles, an alternative is to use a graphical device for
comparing a distribution to another known as a quantile-quantile, or \(q q\) , plot.
 Figure 1.5 illustrates a \(q q\) plot for the bodily injury data using the normal curve
 as a reference distribution. For each point, the vertical axis gives the quantile using
Points in a qq plot the sample distribution. The horizontal axis gives the corresponding quantity
close to a straight line using the normal curve. For example, earlier we considered the 75th percentile
suggest agreement point. This point appears as (1.168, 0.675) on the graph. To interpret a \(q \, q\) plot, if
 between the sample the quantile points lie along the superimposed line, then the sample and the normal
 and the reference reference distribution have the same shape. (This line is defined by connecting
 distributions.
the 75th and 25th percentiles.)
\clearpage

\end{document}
% Compile this UTF-8 document with XeLaTeX or LuaLaTeX.
\documentclass[11pt]{article}
\usepackage{fontspec}
\usepackage{amsmath,amssymb,mathtools,bm}
\usepackage[margin=1in]{geometry}
\setlength{\parindent}{0pt}
\setlength{\parskip}{0.5em}
\begin{document}
\section*{1.3 Power Transforms}
{\small Source: Regression Modeling with Actuarial and Financial Applications (Frees)/Regression Modeling with Actuarial and Financial Applications.pdf\par}
\medskip
1.3 Power Transforms 7
Density Figure 1.6
Distribution of bodily
0.20 injury claims.
Observations are in
0.15 (thousands of)
dollars, with the
0.10 lamges observation
0.05
0.00
0 5 10 15 20
Claims
In Figure 1.5, the small sample percentiles are consistently smaller than the
 corresponding values from the standard normal, indicating that the distribution is
 skewed to the left. The difference in values at the ends of the distribution are due
 to the outliers noted earlier that can also be interpreted as the sample distribution
 having larger tails than the normal reference distribution.
1.3 Power Transforms
In the Section 1.2 example, we considered claims without justifying the use
 of the logarithmic scaling. When analyzing variables such as assets of firms,
wages of individuals, and housing prices of households in business and economic
applications, it is common to consider logarithmic units instead of the original
 units. A log transform retains the original ordering (e.g., large wages remain large
 on the log wage scale) but serves to pull in extreme values of the distribution.
To illustrate, Figure 1.6 shows the bodily injury claims distribution in (thou-
sands of) dollars. To graph the data meaningfully, the largest observation
(\$50,000) was removed prior to making this plot. Even with this observation
removed, Figure 1.6 shows that the distribution is heavily lopsided to the right,
 with several large values of claims appearing. A right-skewed
Distributions that are lopsided in one direction or the other are known as distribution has long.
skewed. Figure 1.6 is an example of a distribution skewed to the right, or positivelytails on the right and
a concentration of
skewed. Here, the tail of the distribution on the right is longer, and there is a greater mass on the left. Many
concentration of mass to the left. In contrast, a left-skewed, or negatively skewed, insurance claims
distribution has a longer tail on the left and a greater concentration of mass toskewedtions areright
 the right. Many insurance claims distributions are right skewed (see Klugman,
Panjer, and Willmot, 2008, for extensive discussions). As we saw in Figures 1.4
and 1.5, a logarithmic transformation yields a distribution that is only mildly
 skewed to the left.
Logarithmic transformations are used extensively in applied statistics work.
 One advantage is that they serve to symmetrize distributions that are skewed. More
 generally, we consider power transforms, also known as the Box-Cox family of 
\clearpage

\end{document}
% Compile this UTF-8 document with XeLaTeX or LuaLaTeX.
\documentclass[11pt]{article}
\usepackage{fontspec}
\usepackage{amsmath,amssymb,mathtools,bm}
\usepackage[margin=1in]{geometry}
\setlength{\parindent}{0pt}
\setlength{\parskip}{0.5em}
\begin{document}
\section*{1.4 Sampling and the Role of Normality}
{\small Source: Regression Modeling with Actuarial and Financial Applications (Frees)/Regression Modeling with Actuarial and Financial Applications.pdf\par}
\medskip
8 Regression and the Normal Distribution
Figure 1.7 500 simu- Density Density
 lated observations 0.00005 0.008
from a chi-square 0.00004 0.006
distribution. The 0.00003 0.004
upper-left panel is
 based on the original 0.00002
distribution. The 0.00001 0.002
upper-right panel 0 0.000
corresponds to the
 square root transform, 0 50,000 100,000 200,000 0 100 200 300 400
 the lower left to the Density y Density Square root of \(y\)
log transform, and the 0.4 100
lower right to the
 negative reciprocal 0.3 80
transform. 60
0.2 40
0.1 20
0.0 r 0
2 4 6 8 10 12 0.15 0.10 0.05 0.00
 Logarithmic \(y\) Negative reciprocal of \(y\)
 transforms. In this family of transforms, in lieu of using the response \(y\) , we use
 a transformed, or rescaled version, \(y^{\lambda}\) . Here, the power \(\lambda\) (lambda, a Greek letter
\(\epsilon\varsigma\mathrm{e l "}\) is a number that may be user specified. Typical values of \(\lambda\) that are used
in practice are \(\lambda=1 , \, 1 / 2 , \, 0 , \, \mathrm{o r} \, \,-1\) . When we use \(\lambda=0\) , we mean \(\operatorname{l n} ( y )\) , that
 is, the natural logarithmic transform. More formally, the Box-Cox family can be
expressed as
\[
y^{( \lambda)}=\begin{cases} {\frac{y^{\lambda}-1} {\lambda}} & {\lambda\neq0} \\ {\operatorname{l n} ( y )} & {\lambda=0} \\ \end{cases} \, .
\]
As we will see, because regression estimates are not affected by location and scale
shifts, in practice, we do not need to subtract \(1\) or divide by \(\lambda\) when rescaling the
response. The advantage of the foregoing expression is that, if we let \(\lambda\) approach
\(\; \; \; \; \; \; \; \; \; \; \; \; \; \; \; \; \; \; \; \; \; \; \; \; \; \; \; \; \; \; \; \; \; \; \; \; \; \; \; \; \; \; \; \; \; \; \; \; \; \; \; \; \; \; \; \; \; \; \; \; \; \; \; \; \; \; \; \; \; \; \; \; \; \; \; \; \; \; \; \; \; \; \; \; \; \; \; \; \; \; \; \; \; \; \; \; \; \; \; \; \; \; \; \; \; \; \; \; \; \; \; \; \; \; \; \; \; \; \; \; \; \; \; \; \; \; \; \; \; \; \; \; \; \; \; \; \; \; \; \; \; \; \; \; \; \; \; \; \; \; \; 0 \; \; \; \; \; 0\) then \(y^{( \lambda)}\) approaches \(\operatorname{l n} ( y )\) , from some straightforward calculus arguments.
To illustrate the usefulness of transformations, we simulated 500 observations
from a chi-square distribution with two degrees of freedom. Appendix A3.2 intro-
duces this distribution (which we will encounter again in studying the behavior
of test statistics). The upper-left panel of Figure 1.7 shows that the original dis-
tribution is heavily skewed to the right. The other panels in Figure 1.7 show the
data rescaled using the square root, logarithmic, and negative reciprocal trans-
 formations. The logarithmic transformation, in the lower-left panel, provides the
 best approximation to symmetry for this example. The negative reciprocal trans-
formation is based on \(\lambda=-1\) and then multiplying the rescaled observations by
\(- 1\) , so that large observations remain large.
1.4 Sampling and the Role of Normality
A statistic is a summary measure of data, such as a mean, median, or per-
centile. Collections of statistics are very useful for analysts, decision makers, and
\clearpage
1.4 Sampling and the Role of Normality
everyday consumers for understanding massive amounts of data that represent
complex situations. To this point, our focus has been on introducing sensible tech-
niques to summarize variables; techniques that will be used repeatedly thought
this text. However, the true usefulness of the discipline of statistics is its ability to
 say something about the unknown, not merely to summarize information already
available. To this end, we need to make some fairly formal assumptions about
the manner in which the data are observed. As a science, a strong feature of
the discipline of statistics is the ability to critique these assumptions and offer
improved alternatives in specific situations.
 It is customary to assume that the data are drawn from a larger population that A statistic is a
we are interested in describing. The process of drawing the data is known as the summary measure of a
 sampling, or data generating, process. We denote this sample as \(\{y_{1} , \, \dots, \, y_{n} \}\) Soample Satistic as
that wemay critieandmodytheesmpling assmptions,we list them hereisci line cne.
used to infer behavior
in detail: about a larger
 population from \(a\)
 sample.
Basic Sampling Assumptions
1.E \(y_{i}=\mu\)
 2. Var \(y_{i}=\sigma^{2}\)
3. \(\{y_{i} \}\) are independent.
4. \(\{y_{i} \}\) are normally distributed.
In this basic setup, \(\mu{}\) and \(\sigma^{2}\) serve as parameters that describe the location and Assumption 4 is not
 on the basis of statistics such as \(\overline{{y}}\) scale of the parent population. The goal is to infer something sensible about themrequired for many
and \(s_{y}^{2}\) For the third assumption, we assume statistical inference
 independence among the draws. In a sampling scheme, this may be uaranteedprocedures becose
central limit theorems
by taking a simple random sample from a population. The fourth assumption provide approximate
 is not required for many statistical inference procedures because central limit normality for many
theorems provide approximate normality for many statistics of interest. However,statisrics of interest.
 a formal justification of some statistics, such as \(t\) -statistics, requires this additional
assumption.
Section 1.9 provides an explicit statement of one version of the central limit
 theorem, giving conditions in which \(\overline{{y}}\) is approximately normally distributed. This
section also discusses a related result, known as an Edgeworth approximation,
 that shows that the quality of the normal approximation is better for symmetric
 parent populations when compared to skewed distributions. Linear regression is
 far we have focused only on the simple arithmetic average \(\overline{{y}}\) Wleeeg
How does this discussion apply to the study of regression analysis? After all, so
 ters, we will emphasize that linear regression is the study of weighted averages;
 specifically, many regression coefficients can be expressed as weighted averages
with appropriately chosen weights. Central limit and Edgeworth approximation
 theorems are available for weighted averages -- these results will ensure approx-
 imate normality of regression coefficients. To use normal curve approximations
 in a regression context, we will often transform variables to achieve approximate
 symmetry.
\clearpage

\end{document}
% Compile this UTF-8 document with XeLaTeX or LuaLaTeX.
\documentclass[11pt]{article}
\usepackage{fontspec}
\usepackage{amsmath,amssymb,mathtools,bm}
\usepackage[margin=1in]{geometry}
\setlength{\parindent}{0pt}
\setlength{\parskip}{0.5em}
\begin{document}
\section*{1.5 Regression and Sampling Designs}
{\small Source: Regression Modeling with Actuarial and Financial Applications (Frees)/Regression Modeling with Actuarial and Financial Applications.pdf\par}
\medskip
10 Regression and the Normal Distribution
Table 1.3
Terminology for y-Variable \(x\) -Variable
Regression Variables
Outcome of interest Explanatory variable
Dependent variable Independent variable
Endogenous variable Exogenous variable
Response Treatment
Regressand Regressor
Left-hand-side variable Right-hand-side variable
 Explained variable Predictor variable
Output Input
1.5 Regression and Sampling Designs
We often transform Approximating normality is an important issue in practical applications of linear
variables to achieve regression. Parts \(\mathrm{I}\) and \(\amalg\) of this book focus on linear regression, where we
 approximate will learn basic regression concepts and sampling design. Part III will focus on
 symmetry to use
normal curve nonlinear regression, involving binary, count, and fat-tailed responses, where
 approximations in a the normal is not the most helpful reference distribution. Ideas concerning basic
regression context. concepts and design are also used in the nonlinear setting.
In regression analysis, we focus on one measurement of interest: the dependent
variable. Other measurements are used as explanatory variables. A goal is to
compare differences in the dependent variable in terms of differences in the
explanatory variables. As noted in Section 1.1, regression is used extensively in
many scientific fields. Table 1.3 lists alternative terms that you may encounter as
 you read regression applications.
In the latter part of the nineteenth century and early part of the twentieth cen-
 tury, statistics was beginning to make an important impact on the development of
experimental science. Experimental sciences often use designed studies, where
the data are under the control of an analyst. Designed studies are performed in
laboratory settings, where there are tight physical restrictions on every variable
 that a researcher thinks may be important. Designed studies also occur in larger
field experiments, where the mechanisms for control are different than in lab-
oratory settings. Agriculture and medicine use designed studies. Data from a
In designed studies, designed study are said to be experimental data.
 the data are under the To illustrate, a classic example is to consider the yield of a crop such as corn,
control of an analyst. where each of several parcels of land (the observations) are assigned various
Dad from saidesinedlevels of fertilizer. The goal is ascertain the effect of fertilizer (the explanatory
 study are said to be
experimental data. variable) on the corn yield (the response variable). Although researchers attempt
 to make parcels of land as much alike as possible, differences inevitably arise.
Agricultural researchers use randomization techniques to assign different levels
of fertilizer to each parcel of land. In this way, analysts can explain the variation
in corn yields in terms of the variation of fertilizer levels. Through the use of
randomization techniques, researchers using designed studies can infer that the
treatment has a causal effect on the response. Chapter \(6\) discusses causality
 further.
\clearpage
1.5 Regression and Sampling Designs 11
Example: Rand Health Insurance Experiment. How are medical-care expen-
ditures related to the demand for insurance? Many studies have established a
 positive relation between the amount spent on medical care and the demand for
health insurance. Those in poor health anticipate using more medical services
 than similarly positioned people in good or fair health and seek higher levels of
 health insurance to compensate for the anticipated expenditures. They obtain this
 additional health insurance by (i) selecting a more generous health insurance plan
from an employer, (ii) choosing an employer with a more generous health insur-
 ance plan, or (iii) paying more for individual health insurance. Thus, it is difficult 
 to disentangle the cause-and-effect relationship of medical-care expenditures and
 the availability of health insurance.
 A study reported by Manning et al. (1987) sought to answer this question using
 a carefully designed experiment. In this study, enrolled households from six cities,
 between November 1974 and February 1977, were randomly assigned to one of
14 different insurance plans. These plans varied by the cost-sharing elements,
the co-insurance rate (the percentage paid on out-of-pocket expenditures, which
varied by \(0 \surd\sigma\) , 25\%, 50\%, and 95\%) as well as the deductible (5\%, 10\%, or
15\% of family income, up to a maximum of \$1,000). Thus, there was a random
assignment to levels of the treatment, the amount of health insurance. The study
 found that more favorable plans resulted in greater total expenditures, even after
controlling for participants' health status.
For actuarial science and other social sciences, designed studies are the excep-
 tion rather than the rule. For example, if we want to study the effects of smoking
on mortality, it is highly unlikely that we could get study participants to agree
 to be randomly assigned to smoker and nonsmoker groups for several years
just so that we could observe their mortality patterns! As with the Section 1.1
Galton study, social science researchers generally work with observational data.
Observational data are not under control of the analyst.
With observational data, we cannot infer causal relationships, but we can
readily introduce measures of association. To illustrate, in the Galton data, it
is apparent that " tall"parents are likely to have " tall"children, and conversely
" short"parents are likely to have " short" children. Chapter \(2\) will introduce a
 correlation and other measures of association. However, we cannot infer causality
from the data. For example, there may be another variable, such as family diet,
 that is related to both variables. Good diet in the family could be associated with
tall heights of parents and adult children, whereas poor diet might stifle, growth.
If this were the case, we would call family diet a confounding variable. With statistical
In designed experiments such as the Rand HealthInsurance Experiment, wecomrar we seion \(\mathbf{x}\)
can control for the effects of variables such as health status through random compare a \({\bf y}\)
assignment methods. In observational studies,we use statistical control ratherecetroling for"the
effects of other
than experimental control. To illustrate, for the Galton data, we might split our explanatory variables.
observations into two groups, one for good family diet and one for poor family
\clearpage

\end{document}
% Compile this UTF-8 document with XeLaTeX or LuaLaTeX.
\documentclass[11pt]{article}
\usepackage{fontspec}
\usepackage{amsmath,amssymb,mathtools,bm}
\usepackage[margin=1in]{geometry}
\setlength{\parindent}{0pt}
\setlength{\parskip}{0.5em}
\begin{document}
\section*{1.6 Actuarial Applications of Regression}
{\small Source: Regression Modeling with Actuarial and Financial Applications (Frees)/Regression Modeling with Actuarial and Financial Applications.pdf\par}
\medskip
12 Regression and the Normal Distribution
 diet, and examine the relationship between parents' and children's height for each
subgroup. This is the essence of the regression method, to compare a \(y\) and an \(x\) :
" controllingfor" the effects of other explanatory variables.
Of course, to use statistical control and regression methods, one must record
family diet and any other measures of height that may confound the effects
 of parents' height on the height of their adult child. The difficulty in designing
 studies is trying to imagine all of the variables that could possibly affect a response
variable, an impossible task in most social science problems of interest. To give
some guidance on when enough is enough, Chapter \(6\) discusses measures of an
explanatory variable's importance and its impact on model selection.
1.6 Actuarial Applications of Regression
This book introduces the statistical method of regression analysis. The introduc-
 tion is organized around the traditional triad of statistical inference:
 Hypothesis testing
: Estimation
: Prediction
Further, this book shows how this methodology can be used in applications that
are likely to be of interest to actuaries and to other risk analysts. As such, it is
 helpful to begin with the three traditional areas of actuarial applications:
. Pricing
. Reserving
: Solvency testing
Pricing and Adverse Selection
Regression analysis can be used to determine insurance prices for many lines
of business. For example, in private passenger automobile insurance, expected
claims vary by the insured's sex, age, location (city versus rural), vehicle purpose
(work or pleasure), and a host of other explanatory variables. Regression can
be used to identify the variables that are important determinants of expected
claims.
In competitive markets, insurance companies do not use the same price for all
 insureds. If they did, " goodrisks," those with lower-than-average expected claims,
would overpay and leave the company. In contrast, "- badrisks," those with higher-
than-average expected claims, would remain with the company. If the company
continued this flat-rate pricing policy, premiums would rise (to compensate for
claims by the increasing share of bad risks) and market share would dwindle
 as the company loses good risks. This problem is known as adverse selection.
Using an appropriate set of explanatory variables, classification systems can be
 developed so that each insured pays his or her fair share.
\clearpage

\end{document}
% Compile this UTF-8 document with XeLaTeX or LuaLaTeX.
\documentclass[11pt]{article}
\usepackage{fontspec}
\usepackage{amsmath,amssymb,mathtools,bm}
\usepackage[margin=1in]{geometry}
\setlength{\parindent}{0pt}
\setlength{\parskip}{0.5em}
\begin{document}
\section*{1.7 Further Reading and References}
{\small Source: Regression Modeling with Actuarial and Financial Applications (Frees)/Regression Modeling with Actuarial and Financial Applications.pdf\par}
\medskip
Chapter References 13
Reserving and Solvency Testing
Both reserving and solvency testing are concerned with predicting whether lia-
bilities associated with a group of policies will exceed the capital devoted to
meeting obligations arising from the policies. Reserving involves determining
the appropriate amount of capital to meet these obligations. Solvency testing is
 about assessing the adequacy of capital to fund the obligations for a block of
 business. In some practice areas, regression can be used to forecast future obli-
gations to help determine reserves (see, e.g., Chapter 19). Regression can also
be used to compare characteristics of healthy and financially distressed firms for
solvency testing (see, e.g., Chapter \(1 4\) .
 Other Risk Management Applications
Regression analysis is a quantitative tool that can be applied in a broad variety of
 business problems, not just in the traditional areas of pricing, reserving, and sol-
vency testing. By becoming familiar with regression analysis, actuaries will have
 another quantitative skill that can be brought to bear on general problems involv-
 ing the financial security of people, companies, and governmental organizations.
To help you develop insights, this book provides many examples of potential
 nonactuarial applications through featured vignettes labeled as " \&amples" and
 illustrative datasets.
To help understand potential regression applications, start by reviewing the
 several datasets featured in the Chapter 1 Exercises. Even if you do not complete
the exercises to strengthen your data summary skills (which requires the use of
a computer), a review of the problem descriptions will help you become more
familiar with types of applications in which an actuary might use regression
 techniques.
1.7 Further Reading and References
This book introduces regression and time series tools that are most relevant to
actuaries and other financial risk analysts. Fortunately, there are other sources
that provide excellent introductions to these statistical topics (although not from
 a risk management viewpoint). Particularly for analysts that intend to specialize
in statistics, it is helpful to get another perspective. For regression, I recommend
Weisburg (2005) and Faraway (2005). For time series, Diebold (2004) is a good
source. Moreover, Klugman, Panjer, and Willmot (2008) provide a good intro-
 duction to actuarial applications of statistics; this book is intended to complement
 the Klugman et al. book by focusing on regression and time series methods.
 Chapter References
Beard, Robert E., Teivo Pentikainen, and Erkki Pesonen (1984). Risk Theory: The Stochastic
Basis of Insurance, 3rd ed. Chapman \& Hall, New York.
Diebold, Francis X. (2004). Elements of Forecasting, 3rd ed. Thomson, South-Western, Mason,
Ohio.
\clearpage

\end{document}
% Compile this UTF-8 document with XeLaTeX or LuaLaTeX.
\documentclass[11pt]{article}
\usepackage{fontspec}
\usepackage{amsmath,amssymb,mathtools,bm}
\usepackage[margin=1in]{geometry}
\setlength{\parindent}{0pt}
\setlength{\parskip}{0.5em}
\begin{document}
\section*{1.8 Exercises}
{\small Source: Regression Modeling with Actuarial and Financial Applications (Frees)/Regression Modeling with Actuarial and Financial Applications.pdf\par}
\medskip
14 Regression and the Normal Distribution
 Faraway, Julian J. (2005). Linear Models in R. Chapman \& Hall/CRC, New York.
Hogg, Robert V. (1972). On statistical education. American Statistician 26, 8-- 11.
 Klugman, Stuart A., Harry H. Panjer, and Gordon E. Willmot (2008). Loss Models: From Data
to Decisions. John Wiley \(\&\) Sons, Hoboken, New Jersey.
Manning, Willard G., Joseph P. Newhouse, Naihua Duan, Emmett B. Keeler, Arleen Leibowitz,
 and M. Susan Marquis (1987). Health insurance and the demand for medical care: Evidence
 from a randomized experiment. American Economic Review 77, no. 3, 251-- 277.
Rempala, Grzegorz A., and Richard A. Derrig (2005). Modeling hidden exposures in claim
 severity via the EM algorithm. North American Actuarial Journal 9, no. 2, 108-- 128.
Rosenberg, Marjorie A., Edward W. Frees, Jia Feng Sun, Paul Johnson Jr., and James \(\mathbf{M}\) .
Robinson (2007). Predictive modeling with longitudinal data: A case study of Wisconsin
 nursing homes. North American Actuarial Journal 11, no. 3, 54- 69.
 Singer, Judith \(\mathrm{D}\) , and J. B. Willett (1990). Improving the teaching of applied statistics: Putting
 the data back into data analysis. American Statistician 44, 223-- 230.
Stigler, Steven M. (1986). The History of Statistics: The Measurement of Uncertainty before
1900. Belknap Press, Harvard University Press, Cambridge, Massachusetts.
Weisberg, Sanford (2005). Applied Linear Regression, 3rd ed. John Wiley \& Sons, New York.
@ EMPIRICAL 1.8 Exercises
Filename is 1.1. MEPS Health Expenditures. This exercise considers data from the Med-
"HealthExpend" ical Expenditure Panel Survey (MEPS), conducted by the U.S. Agency
 of Health Research and Quality. MEPS is a probability survey that pro-
vides nationally representative estimates of health-care use, expenditures,
sources of payment, and insurance coverage for the U.S. civilian popu-
lation. This survey collects detailed information on individuals of each
medical-care episode by type of services, including physician office visits,
hospital emergency room visits, hospital outpatient visits, hospital inpatient
stays, all other medical provider visits, and use of prescribed medicines.
This detailed information allows one to develop models of health-care
use to predict future expenditures. You can learn more about MEPS at
http://www.meps.ahrq.gov/mepsweb/.
We consider MEPS data from panels 7 and \(8\) of 2003 that consists of
18,735 individuals between ages 18 and 65. From this sample, we took a
random sample of 2,000 individuals that appear in the file " HealthExpend.
 From this sample, there are 157 individuals that had positive inpatient expen-
ditures. There are also 1,352 that had positive outpatient expenditures. We
 will analyze these two samples separately.
Our dependent variables consist of amounts of expenditures for inpatient
(EXPENDIP) and outpatient (EXPENDOP) visits. For MEPS, outpatient
events include hospital outpatient department visits, office-based provider
visits, and emergency room visits excluding dental services. (Dental ser-
vices, compared to other types of health-care services, are more predictable
and occur on a more regular basis.) Hospital stays with the same date
of admission and discharge, known as zero-night stays, were included in
outpatient counts and expenditures. (Payments associated with emergency
room visits that immediately preceded an inpatient stay were included in
\clearpage
1.8 Exercises 15
the inpatient expenditures. Prescribed medicines that can be linked to hos-
pital admissions were included in inpatient expenditures, not in outpatient
utilization.)
 Part 1: Use only the 157 individuals who had positive inpatient expen-
ditures and do the following analysis:
a.Compute descriptive statistics for inpatient (EXPENDIP) expenditures.
\(\mathrm{a}\) (i) What is the typical (mean and median) expenditure?
\(\mathrm{a}\) (ii) How does the standard deviation compare to the mean? Do the data
appear to be skewed?
b. Compute a box plot, histogram, and a (normal) \(q q\) plot for EXPENDIP.
Comment on the shape of the distribution.
c.Transformations.
c(i) Take a square root transform of inpatient expenditures. Summarize
 the resulting distribution using a histogram and a \(q \, q\) plot. Does it
appear to be approximately normally distributed?
c(ii) Take a (natural) logarithmic transformation of inpatient expendi-
tures. Summarize the resulting distribution using a histogram and a
\(q q\) plot. Does it appear to be approximately normally distributed?
Part \({\bf2}\) : Use only the 1,352 individuals who had positive outpatient
expenditures.
d. Repeat part (a) and compute histograms for expenditures and loga-
rithmic expenditures. Comment on the approximate normality for each
histogram. @ EMPIRICAL
1.2. Nursing Home Utilization. This exercise considers nursing home data pro- Filename is
vided by the Wisconsin Department of Health and Family Services (DHFS)."WiscNursing Home"
The State of Wisconsin Medicaid program funds nursing home care for
individuals qualifying on the basis of need and financial status. As part
of the conditions for participation, Medicaid-certified nursing homes must
file an annual cost report to DHFS, summarizing the volume and cost of
care provided to all of its residents, Medicaid funded and otherwise. These
cost reports are audited by DHFS staff and form the basis for facility-
specific Medicaid daily payment rates for subsequent periods. The data are
publicly available; see http://dhfs.wisconsin.gov/provider/prev-yrs-reports-
nh.htm for more information.
 The DHFS is interested in predictive techniques that provide reliable uti-
lization forecasts to update its Medicaid-funding rate schedule of nursing
facilities. In this assignment, we consider the data in the file " WscNurs-
ingHome" in cost-report years 2000 and 2001. There are 362 facilities in
 2000 and 355 facilities in 2001. Typically, utilization of nursing home care
is measured in patient days (i.e., the number of days each patient was in the
facility, summed over all patients). For this exercise, we define the outcome
variable to be total patient years (TPY), the number of total patient days
in the cost-reporting period divided by number of facility operating days
in the cost-reporting period (see Rosenberg et al., 2007, Appendix 1, for
further discussion of this choice). The number of beds (NUMBED) and
\clearpage
16 Regression and the Normal Distribution
square footage (SQRFOOT) of the nursing home both measure the size of
the facility. Not surprisingly, these continuous variables will be important
 predictors of TPY.
Part 1: Use cost-report year 2000 data, and do the following analysis:
a. Compute descriptive statistics for TPY, NUMBED, and SQRFOOT.
b. Summarize the distribution of TPY using a histogram and a \(q \, q\) plot.
Does it appear to be approximately normally distributed?
c. Transformations. Take a (natural) logarithmic transformation of TPY
(LOGTPY). Summarize the resulting distribution using a histogram
 and a \(q \, q\) plot. Does it appear to be approximately normally distri-
buted?
@ EMPIRICAL Part 2: Use cost-report year 2001 data and repeat parts (a) and (c).
Filename is 1.3. Automobile Insurance Claims. As an actuarial analyst, you are working
"AutoClaims" with a large insurance company to help it understand claims distribution for
 private passenger automobile policies. You have available claims data for a
recent year, consisting of
: STATE CODE: codes 01 through 17 used, with each code randomly
 assigned to an actual individual state
: CLASS: rating class of operator, based on age, sex, marital status, and
 use of vehicle
: GENDER: operator sex AGE: operator age
- PAID: amount paid to settle and close a claim.
You are focusing on older drivers, 50 and older, for which there are
\(n \,=6 , \! 7 7 3\) claims available.
Examine the histogram of the amount PAID and comment on the symme-
try. Create a new variable, the (natural) logarithmic claims paid, LNPAID.
Create a histogram and a \(q q\) plot of LNPAID. Comment on the symmetry
@ EMPIRICAL of this variable.
Filename is 1.4. Hospital Costs. Suppose that you are an employee benefits actuary working
"HospitalCosts" with a medium-sized company in Wisconsin. This company is considering
offering, for the first time in its industry, hospital insurance coverage to
dependent children of employees. You have access to company records
and so have available the number, age, and sex of the dependent children
but have no other information about hospital costs from the company. In
 particular, no firm in this industry has offered this coverage, and so you have
little historical industry experience on which you can forecast expected
 claims.
You gather data from the Nationwide Inpatient Sample of the Healthcare
Cost and Utilization Project (NIS-HCUP), a nationwide survey of hospital
costs conducted by the U.S. Agency for Healthcare Research and Quality
(AHRQ). You restrict consideration to Wisconsin hospitals and analyze
\(\bf a\) random sample of \(n=5 0 0\) claims from 2003 data. Although the data
come from hospital records, they are organized by individual discharge, and
so you have information about the age and sex of the patient discharged.
Specifically, you consider patients aged 0-- 17 years. In a separate project,
\clearpage
1.8 Exercises 17
you will consider the frequency of hospitalization. For this project, the goal
is to model the severity of hospital charges, by age and sex.
a.Examine the distribution of the dependent variable, TOTCHG. Do this
 by making a histogram and then a \(q \, q\) plot, comparing the empirical to
 a normal distribution.
 b.Take a natural log transformation and call the new variable LNTOTCHG.
Examine the distribution of this transformed variable. To visualize the
logarithmic relationship, plot LNTOTCHG versus TOTCHG. @ EMPIRICAL
1.5. Automobile Injury Insurance Claims. We consider automobile injury Filename is
claims data using data from the Insurance Research Council (IRC), a divi- "AutoBI"
 sion of the American Institute for Chartered Property Casualty Underwriters
and the Insurance Institute of America. The data, collected in 2002, con-
 tain information on demographic information about the claimant, attorney
involvement, and economic loss (LOSS, in thousands), among other vari-
ables. We consider here a sample of \(n \,=\, 1 , 3 4 0\) losses from a single state.
The full 2002 study contains more than 70,000 closed claims based on data
from 32 insurers. The IRC conducted similar studies in 1977, 1987, 1992,
 and 1997.
a.Compute descriptive statistics for the total economic loss (LOSS). What
is the typical loss?
b. Compute a histogram and (normal) \(q \, q\) plot for LOSS. Comment on the
 shape of the distribution.
c. Partition the dataset into two subsamples, one corresponding to those
claims that involved an ATTORNEY \((=1 )\) and the other to those in
which an ATTORNEY was not involved \((=2 )\) .
\(\mathrm{c ( i )}\) For each subsample, compute the typical loss. Does there appear
to be a difference in the typical losses by attorney involvement?
c(ii) To compare the distributions, compute a box plot by level of
 attorney involvement.
c(iii) For each subsample, compute a histogram and \(q q\) plot. Compare
the two distributions. @ EMPIRICAL
1.6. Insurance Company Expenses. Like other businesses, insurance compa- Filename is
nies seek to minimize expenses associated with doing business to enhance "NAICExpense"
profitability. To study expenses, this exercise examines a random sample of
 500 insurance companies from the National Association of Insurance Com-
missioners' (NAIC) database of more than 3,000 companies. The NAIC
maintains one of the world's largest insurance regulatory databases; we
 consider here data that are based on 2005 annual reports for all the property
 and casualty insurance companies in the United States. The annual reports
 are financial statements that use statutory accounting principles.
Specifically, our dependent variable is EXPENSES, the nonclaim ex-
penses for a company. Although not needed for this exercise, nonclaim
 expenses are based on three components: unallocated loss adjustment, under-
writing, and investment expenses. The unallocated loss adjustment expense
is the expense not directly attributable to a claim but indirectly associated
\clearpage

\end{document}
% Compile this UTF-8 document with XeLaTeX or LuaLaTeX.
\documentclass[11pt]{article}
\usepackage{fontspec}
\usepackage{amsmath,amssymb,mathtools,bm}
\usepackage[margin=1in]{geometry}
\setlength{\parindent}{0pt}
\setlength{\parskip}{0.5em}
\begin{document}
\section*{1.9 Technical Supplement - Central Limit Theorem}
{\small Source: Regression Modeling with Actuarial and Financial Applications (Frees)/Regression Modeling with Actuarial and Financial Applications.pdf\par}
\medskip
18 Regression and the Normal Distribution
with settling claims; it includes items such as the salaries of claims adjusters,
legal fees, court costs, expert witnesses, and investigation costs. Underwrit-
ing expenses consist of policy acquisition costs, such as commissions, as
 well as the portion of administrative, general, and other expenses attributable
to underwriting operations. Investment expense are those expenses related
to investment activities of the insurer.
a. Examine the distribution of the dependent variable, EXPENSES. Do
 this by making a histogram and then a \(q \, q\) plot, comparing the empirical
 to a normal distribution.
 b. Take a natural log transformation and examine the distribution of this
 transformed variable. Has the transformation helped to symmetrize the
@ EMPIRICAL distribution?
Filename is 1.7. National Life Expectancies. Who is doing health care right? Health-care
"UNLifeExpectancy" decisions are made at the individual, corporate, and government levels.
Virtually every person, corporation, and government has a different perspec-
tive on health care; these result in a wide variety of systems for managing
health care. Comparing different health-care systems help us learn about
 approaches other than our own, which in turn help us make better decisions
in designing improved systems.
Here, we consider health-care systems from \(n \,=\, 1 8 5\) countries through-
 out the world. As a measure of the quality of care, we use LIFEEXP, the life
expectancy at birth. This dependent variable and several explanatory vari-
 ables are listed in Table 1.4. From this table, you will note that although there
 are 185 countries consider in this study, not all countries provided informa-
 tion for each variable. Data not available are noted under the column " Num
Miss." The data are from the UN Human Development Report.
a.Examine the distribution of the dependent variable, LIFEEXP. Do this
by making a histogram and then a \(q \, q\) plot, comparing the empirical to
a normal distribution.
 b. Take a natural log transformation and examine the distribution of this
 transformed variable. Has the transformation helped to symmetrize the
distribution?
 1.9 Technical Supplement -- Central Limit Theorem
Central limit theorems form the basis for much of the statistical inference used
in regression analysis. Thus, it is helpful to provide an explicit statement of one
version of the central limit theorem.
 Central Limit Theorem. Suppose that \(y_{1} , \dots, \, y_{n}\) are independently distributed
 with mean \(\mu\) , finite variance \(\sigma^{2}\) and \(\mathrm{E} | y |^{3}\) is finite. Then,
\[
\underset{n \rightarrow\infty} {\operatorname* {l i m}} \operatorname* {P r} \left( \frac{\sqrt{n}} {\sigma} ( \overline{{y}}-\mu) \leq x \right)=\Phi\left( x \right)
\]
 for each \(x\) , where \(\Phi\left( . \right)\) is the standard normal distribution function.
\clearpage
1.9 Technical Supplement -- Central Limit Theorem 19
 Table 1.4 Life
Num Standard Expectancy,
 Variable Description Miss Mean Median Deviation Minimum Maximum Economic, and
Demographic
BIRTH Births attended 7 78.2592.00 26.42 6.00 100.00 Characteristics of 185
ATTEND by skilled Countries
health
personnel (\%)
FEMALE Legislators, 87 29.0730.00 11.71 2.00 58.00
BOSS senior
officials, and
managers, \(\mathcal{J}_{o}\)
 female
FERTILITYTotal fertility 4 3.19 2.70 1.71 0.90 7.50
rate, births per
woman
GDP Gross domestic 7 247.55 14.20 1,055.69 0.10 12,416.50
product, in
 billions of
USD
HEALTH 2004 health 5718.01 297.50 1,037.0115.00 6,096.00
EXPEND expenditure
 per capita,
PPP in USD
ILLITERATE Adult illiteracy 14 17.69 10.1019.86 0.20 76.40
rate, \(\mathcal{J}_{o}\) aged
15 and older
PHYSICIAN Physicians, per 3 146.08 107.50 138.552.00 591.00
100,000
people
POP 2005 population, 1 35.36 7.80 131.70 0.10 1,313.00
 in millions
PRIVATE 2004 private 1 2.52 2.40 1.33 0.30 8.50
HEALTH expenditure on
health, \(\mathcal{T}_{o}\) of
GDP
PUBLIC Public 28 4.69 4.60 2.05 0.60 13.40
EDUCATIONexpenditure on
education, \(\gamma_{o}\)
of GDP
RESEAR Researchers in 95 2,034.66 848.00 4,942.93 15.00 45,454.00
CHERS R\&D,per
 million people
SMOKING Prevalence of 88 35.09 32.00 14.40 6.00 68.00
 smoking,
(male) \% of
adults
LIFEEXP Life expectancy 67.05 71.00 11.08 40.50 82.30
at birth, in
years
Source: UN Human Development Report, available at http://hdr.undp.org/en/.
 Note: PPP = Purchasing Power Parity.
\clearpage
20 Regression and the Normal Distribution
Under the assumptions of this theorem, the rescaled distribution of \(\overline{{y}}\) appro-
aches a standard normal as the sample size, \(n\) increases. We interpret this as
 meaning that, for large sample sizes, the distribution of \(\overline{{y}}\) may be approximated
by a normal distribution. Empirical investigations have shown that sample sizes
of \(n=2 5\) to 50 provide adequate approximations for most purposes.
When does the central limit theorem not work well? Some insights are provided
 by another result from mathematical statistics.
Edgeworth Approximation. Suppose that \(y_{1} , \dots, \, y_{n}\) are identically and inde-
pendently distributed with mean \(\mu{}\) , finite variance \(\sigma^{2}\) and \(\mathrm{E} | y |^{3}\) is finite.
Then,
\[
\begin{aligned} {\operatorname* {P r} \left( \frac{\sqrt n} {\sigma} ( \overline{y}-\mu) \leq x \right)} & {{}=\Phi\left( x \right)+\frac1 6 \frac1 {\sqrt{2 \pi}} e^{-x^{2} / 2} \frac{\mathrm{E} \left( y-\mu\right)^{3}} {\sigma^{3} \sqrt n}+\frac{h_{n}} {\sqrt n}} \\ \end{aligned}
\]
for each \(x\) , where \(h_{n} \to0 \mathrm{~ a s ~} n \to\infty\)
This result suggests that the distribution of \(\overline{{y}}\) becomes closer to a normal dis-
tribution as the skewness, \(\mathrm{E} ( \overline{{y}}-\mu)^{3}\) , becomes closer to zero. This is important
in insurance applications because many distributions tend to be skewed. Histori-
cally, analysts used the second term on the right-hand side of the result to provide
\(\mathrm{a}\) " correction"for the normal curve approximation. See, for example, Beard,
Pentikainen and Pesonen (1984) for further discussion of Edgeworth approxi-
mations in actuarial science. An alternative (used in this book) that we saw in
Section 1.3 is to transform the data, thus achieving approximate symmetry. As
suggested by the Edgeworth approximation theorem, if our parent population is
close to symmetric, then the distribution of \(\overline{{y}}\) will be approximately normal.
\clearpage
Part \(\mathbf{I}\)
Linear Regression
\clearpage

\clearpage

\end{document}
